\chapter{Small-$L$, unbounded eta: uniform source weights}
\label{component:p66922bb7f235}
\noindent\textbf{Component }\path{p66922bb7f235}\par
{\small\noindent Source: \path{R13/advances/revision13_c1_small_l/eta_large_bounds.tex}\par}
{\small\noindent SHA256: \path{66922bb7f235f4ea968033105881fd80854a5e65dc0223bf372d0969701ccc0a}\par}
{\small\noindent Inventory use: ES source bounds | C14\_\allowbreak{}ENTIRE\_\allowbreak{}ORIGINAL\_\allowbreak{}PARENT\_\allowbreak{}GRAPH; precise subordinate uses retained below\par}
\medskip
\section{Small-$L$, unbounded eta: uniform source weights}
\label{p66922bb7f235:sec:r13-small-L-eta-large-bounds}

This is a source theorem, not an additional all-tests admission.
Use the original physical C1 source with $x,L=x\zeta>0$
infinitesimal, $\zeta$ noninfinitesimal and $\eta$ unbounded.
Put $\sigma=1/\eta$, $S=\sigma+L$, and retain the original
$0<r,h<1$, $\theta=1-h$, $M>0$, and $D,\beta\ge0$.
The native gap gives $E=1-xh/\sigma\ge1/18$. Set
\[
 e={xh\over\sigma}=1-E,\quad
 P=\sigma P_d,\quad J=1+xh-\sigma(1-r),\quad
 H_0=1+L(1-r),\quad B=P+x(J-\sigma),
\]
\[
 \mathcal H=R H_0+(1-r)xh(2+xr)
 -e\{2+x(1+h+r\theta)\}+Ke^2,
 \quad R=1+xr,\quad K=1+x.                           \tag{ES1}
\]
No $E\sim1$ assumption is made here. The original source formulas
give the two exact identities
\[
 2\theta P\beta=hB-\sigma H_0,\qquad
 2P(JD+x\theta\beta)=\sigma^2\mathcal H.             \tag{ES2}
\]
They are the beta-shared identities with the reciprocal-$L$ factors
cleared. In particular their use does not require bounded $1/L$.

The positive denominator is exactly
\[
 P=\sigma(1+L+2xh+rxL+rx^2h)+KL+x^2h\theta.
                                                               \tag{ES3}
\]
Since $x/L=1/\zeta$ is bounded and $x^2h=x\sigma e$,
all terms beyond $\sigma+L$ in (ES3), divided by $S$, are
infinitesimal. Consequently
\[
 {P\over S}\sim1,\qquad J\sim1,\qquad
 {B\over S}-(1+x/S)\in\Inf,\qquad
 H_0\sim1,\qquad \mathcal H-E^2\in\Inf.             \tag{ES4}
\]
Here $x/S$ is bounded. Both $B/S$ and $\mathcal H$ are
positive bounded units, uniformly through all original openings;
for the latter use the actual bound $E\ge1/18$.

Define nonnegative pure-source sphere coordinates
\[
 d_0^2={2PJD\over\sigma^2\mathcal H},\qquad
 b_0^2={2Px\theta\beta\over\sigma^2\mathcal H}.
\]
Then
\[
 d_0^2+b_0^2=1,\qquad
 D={\sigma^2\mathcal H\over2PJ}d_0^2,\qquad
 x\theta\beta={\sigma^2\mathcal H\over2P}b_0^2.
                                                               \tag{ES5}
\]
Thus $D=O(\sigma^2/S)=O(\sigma)$ and
$x\theta\beta=O(\sigma^2/S)$, including either literal coefficient
boundary. The first identity of (ES2) becomes
\[
 eB=xH_0+\sigma\mathcal H b_0^2.                     \tag{ES6}
\]
Put $\kappa=x/S$ and $\nu=\sigma/S$. Multiplying
$e(B/S)=\kappa H_0+\nu\mathcal H b_0^2$ by $e$ and using
$e\nu=\kappa h$ gives
\[
 e^2(B/S)=\kappa\{eH_0+h\mathcal H b_0^2\}.
                                                               \tag{ES7}
\]
All factors in braces are bounded and nonnegative. Hence
$e^2=O(x/S)$, and (ES6) also gives $x/S=O(e)$.
In particular $\zeta$ unbounded implies $e$ infinitesimal,
since $x/S\le1/\zeta$. The unit-$\zeta$ cases may have
$e$ a positive bounded unit; (SN7) and (SN12) are not imported.

The other original coefficients have the opening-cleared bounds
\[
 {CS(1-r)\over E}\sim\tfrac12,\qquad
 {ASr\over E}\sim\tfrac12,\qquad
 m={K(1+\sigma r)eS\over2M\zeta P},                  \tag{ES8}
\]
\[
 \alpha={(1+\sigma r)\{LR eS+L^2ME\}\over2QP}
       =O(L),\qquad Q=1-Lr\sim1.                   \tag{ES9}
\]
For $A$ use $P_4=(1+xh)E+(1-r)(xh+L)\sim E$.
The bound on $\alpha$ follows because $0<M<R$ is bounded and
$L/S\le1$. Thus $C,A$ are unbounded while $D,\alpha$ are
infinitesimal on the entire stated source domain. No bound on $m$
is asserted through $M=0$, and no bound on beta is asserted through
$\theta=0$.

Finally put $\mu=K(1+\sigma r)S/(2M\zeta P)>0$. The
original $m$ packet has the exact shared-sphere split
\[
 m=m_0+m_1b_0^2,\qquad
 m_0={\mu xH_0\over B},\qquad
 m_1={\mu\sigma\mathcal H\over B}.                  \tag{ES10}
\]
The remaining source opening $M$ is displayed explicitly. These
identities provide valid weights for C14.4 and the boundedness step
missing from its earlier checklist. A full-parent allocation for every
residual infinitesimal-$u$ test is still required. The narrow checker
verifies (ES2), (ES3), (ES6)--(ES10) algebraically; the inequalities
above are the ordered-field argument. No mixed inverse, mixed order,
or source/testing rate comparison is used.
